% answers_short.tex —— 1987 年·广东卷·文史类 · 简短答案速查
% 由 tools/make_short_answers.py 自 answers.tex 的 \gksolhead/\gkanswer 提取，
% 并与 1987/1987_数学真题_广东文_答案要点.md 逐题比对后生成，勿手改（改源 MD 后重跑）。

\begin{gknotice}
\textbf{使用说明}：本节只给最终答案便于快速核对，完整推导见解析版。答案为编写者独立推导，\textbf{非官方标准答案}。
\end{gknotice}

\gkpart{大题}
\begin{itemize}[leftmargin=2.6em]
  \item[\textbf{1.}] (1) A；(2) B；(3) C；(4) D；(5) B；(6) D；(7) A；(8) B；(9) D；(10) C；(11) B；(12) D；(13) B；(14) A；(15) B；(16) D；(17) B；(18) A；(19) B；(20) D；(21) D；(22) C；(23) D；(24) B；(25) D
  \item[\textbf{2.}] (1) $1-2\sqrt2\,i$；(2) $\displaystyle\frac{x^2}{8}+\frac{y^2}{4}=1$；(3) $2$；(4) $64$；(5) $[-1,1]$
  \item[\textbf{3.}] (1) $1$；(2) $a=6,\ b=8,\ c=10$
  \item[\textbf{4.}] (1) 本固定源仅交叉引用印刷页 286；独立推导见解析正文；(2) $\sqrt2\text{ cm}$（独立推导）
  \item[\textbf{5.}] (1) $y_1y_2=-p^2$；(2) $B,C$ 与抛物线顶点三点共线
  \item[\textbf{6.}] (1) $\displaystyle S_{2n}=\frac12(\sin2\theta)\left[1+(-1)^{n+1}\operatorname{tg}^{2n}\theta\right]$；(2) $\displaystyle a_n^2+a_{n-1}a_{n+1}+(-1)^n(a_1)^{2n}=0$
\end{itemize}
